% Standalone certified dossier for the stochastic-localization/Riccati core.
\documentclass[../main.tex]{subfiles}
\begin{document}

% === SOLUTION HEADER (proof provenance) =====================================
% ledger-nodes: lem:survival-implies-kls; lem:matrix-riccati;
%               thm:scalar-riccati; cor:per-direction;
%               cor:tight-window-consumption; lem:pathwise-BL;
%               cor:away-from-zero
% refines: lem:survival-implies-kls; lem:matrix-riccati;
%          thm:scalar-riccati; cor:per-direction;
%          cor:tight-window-consumption; lem:pathwise-BL;
%          cor:away-from-zero
% bounded_by: none (none of the seven ledger nodes carries a bounded_by edge)
% evidence_run: none
% checked_by: agent
% authored_by: /root/kls_core_author
% reviewed_by: /root/kls_bootstrap_author
% review: research/reviews/2026-08-25-kls-core-r2-audit.md
% date: 2026-08-25
% ============================================================================

\section*{Solution: the localization, Riccati, and tight-window core}
\label{sec:sol-kls-localization-riccati-core}

\paragraph{Scope and regularity.}
This dossier proves exactly the seven results listed in the header.  We first work with a
smooth, compactly supported, full-dimensional log-concave probability measure and a smooth
approximation of the cut.  On every bounded time interval all stochastic integrals below are
then true martingales after the usual bounded stopping.  Every identity is local and every
estimate has constants independent of the regularization.  Truncating the measure, smoothing
the density and the cut, putting the approximation in isotropic position, and then using
Fatou's lemma for the nonnegative occupation terms gives the stated general versions.  The
only perimeter statement needed without compact support is the supermartingale inequality
\begin{equation}\label{eq:sol-perimeter-supermartingale}
  \E\mu_t^+(E)\le \mu^+(E).
\end{equation}
It follows directly by integrating, over the reduced boundary of the fixed cut, the
nonnegative local-martingale density of stochastic localization.  Thus no optional-stopping
argument below presupposes unproved uniform integrability at an infinite endpoint.

\subsection*{1. Setup and elementary decompositions}

Let $\mu$ be isotropic and log-concave on $\R^n$.  Eldan localization is the posterior process
\[
  \dd\mu_t(x)=Z_t^{-1}\exp\!\left(c_t\cdot x-\frac t2\abs{x}^2\right)\dd\mu(x),
\]
adapted to an $n$-dimensional Brownian motion $W$.  Write
\[
 a_t=\E_{\mu_t}X,\qquad A_t=\Cov_{\mu_t}(X).
\]
For every integrable test function $\phi$, after localization when necessary,
\begin{equation}\label{eq:sol-loc-martingale}
 \dd\!\int\phi\,\dd\mu_t=\Cov_{\mu_t}(\phi,X)\cdot\dd W_t,
 \qquad
 \dd a_t=A_t\dd W_t,
 \qquad
 \dd A_t=\dd\mathcal T_t-A_t^2\dd t,
\end{equation}
where $\mathcal T$ is a matrix local martingale.  For $t>0$, the posterior is
$t$-uniformly log-concave, hence Brascamp--Lieb gives
\begin{equation}\label{eq:sol-BL-cap}
 A_t\preceq t^{-1}I_n.
\end{equation}

Fix a measurable cut $E$, put $F=E^c$, and, as long as $p_t,q_t>0$, define
\[
 p_t=\mu_t(E),\quad q_t=1-p_t,\quad s_t=p_tq_t,
\]
\[
 m_t^E=\E_{\mu_t}[X\mid E],\quad m_t^F=\E_{\mu_t}[X\mid F],\quad
 \delta_t=m_t^E-m_t^F,
\]
\[
 \Sigma_t^E=\Cov_{\mu_t}(X\mid E),\quad
 \Sigma_t^F=\Cov_{\mu_t}(X\mid F),\quad G_t=\Sigma_t^E-\Sigma_t^F.
\]
Finite-time posteriors have a strictly positive likelihood relative to $\mu$, so an initial
$p_0\in(0,1)$ keeps $p_t,q_t>0$ at every finite time.  Covariance decomposition gives
\begin{equation}\label{eq:sol-cov-decomp}
 A_t=p_t\Sigma_t^E+q_t\Sigma_t^F+s_t\delta_t\delta_t^T.
\end{equation}
Set
\[
 B_t=s_t\delta_t\delta_t^T,
 \quad R_t=A_t-B_t=p_t\Sigma_t^E+q_t\Sigma_t^F\succeq0,
 \quad r_t=\Tr B_t=s_t\abs{\delta_t}^2,
\]
\[
 K_t=G_t+(q_t-p_t)\delta_t\delta_t^T,
 \quad S_t=s_t\norm{G_t}_{\HS}^2,
 \quad D_t=2s_t\delta_t^TA_t\delta_t-r_t^2.
\]
In particular $0\preceq B_t\preceq A_t$, $B_t^2=r_tB_t$, and at time zero
$R_0\preceq A_0=I_n$.

Finally let
\[
 v_t=\int_E(x-a_t)\dd\mu_t(x)=s_t\delta_t.
\]
Then
\begin{equation}\label{eq:sol-p-qv}
 \dd p_t=v_t\cdot\dd W_t,
 \qquad \dd[p]_t=\abs{v_t}^2\dd t=s_t^2\abs{\delta_t}^2\dd t.
\end{equation}

\subsection*{2. Balanced posterior survival gives boundary and KLS}

\begin{lemma}[= Lemma~\ref{lem:survival-implies-kls}]
\label{lem:sol-survival-implies-kls}
Suppose that for some $T_0,c_0,b_0>0$,
\[
 \Prob\{\min(p_{T_0},q_{T_0})\ge b_0\}\ge c_0.
\]
Then
\[
 \mu^+(E)\ge c\,c_0b_0\sqrt{T_0}
\]
for a numerical $c>0$.  Consequently, if $T_0,c_0,b_0$ are universal and the event holds for
every balanced cut of every isotropic log-concave measure, then the KLS Cheeger constant is
bounded below universally.
\end{lemma}

\begin{proof}
The potential of $\mu_{T_0}$ has Hessian bounded below by $T_0I$.  The standard
Bakry--Emery/Bobkov isoperimetric consequence for a $T_0$-uniformly log-concave law is
\[
 \mu_{T_0}^+(E)\ge c\sqrt{T_0}\min(p_{T_0},q_{T_0});
\]
see \cite{BakryGentilLedoux2014,Bobkov1999LogConcave,Milman2009Isoperimetric}.  Therefore
\eqref{eq:sol-perimeter-supermartingale} and the assumed event give
\[
 \mu^+(E)\ge\E\mu_{T_0}^+(E)
 \ge c\sqrt{T_0}\E\min(p_{T_0},q_{T_0})
 \ge c\,c_0b_0\sqrt{T_0}.
\]
The isoperimetric profile of a log-concave measure is concave, symmetric under
$p\leftrightarrow1-p$, and vanishes at the endpoints.  Hence a common positive lower bound for
$I_\mu(1/2)$, equivalently for every balanced cut by approximation of the profile infimum,
gives $I_\mu(p)\ge c'\min(p,1-p)$ for all $p$.  This is the dimension-free Cheeger form of
KLS.
\end{proof}

\subsection*{3. The matrix and scalar Riccati identities}

\begin{lemma}[= Lemma~\ref{lem:matrix-riccati}]
\label{lem:sol-matrix-riccati}
There are matrix local martingales $N,\widetilde N$ such that
\begin{align}
 \dd B_t&=\dd N_t+
 \bigl(s_tG_t^2-A_tB_t-B_tA_t+r_tB_t\bigr)\dd t,
 \label{eq:sol-dB}\\
 \dd R_t&=\dd\widetilde N_t-
 \bigl(R_t^2+s_tG_t^2\bigr)\dd t.
 \label{eq:sol-dR}
\end{align}
\end{lemma}

\begin{proof}
Suppress the time subscript.  Define the unnormalised conditional second moment
\[
 C^E=\int_E(x-a)(x-a)^T\dd\mu_t(x).
\]
Applying \eqref{eq:sol-loc-martingale} to $\int_E x\dd\mu_t$ and using It\^o's product rule
for $pa$ gives
\begin{equation}\label{eq:sol-dv}
 \dd v=(C^E-pA)\dd W-Av\dd t.
\end{equation}
Indeed, the quadratic covariation of $p$ and $a$ is $Av\dd t$.  From
\eqref{eq:sol-cov-decomp}, $m^E-a=q\delta$, and $m^F-a=-p\delta$, one obtains
\[
 C^E-pA=s\bigl(G+(q-p)\delta\delta^T\bigr)=sK.
\]
Thus $\dd v=sK\dd W-Av\dd t$.  Since $s=p(1-p)$ and $\dd p=v\cdot\dd W$,
\begin{equation}\label{eq:sol-ds}
 \dd s=(q-p)v\cdot\dd W-\abs v^2\dd t.
\end{equation}

For completeness, we retain every It\^o correction in $B=vv^T/s$.  Put $\alpha=q-p$.
The drift of $vv^T$ is
\[
 -Avv^T-vv^TA+s^2K^2,
\]
while
\[
 \dd(s^{-1})=-\alpha s^{-2}v\cdot\dd W
 +\left(\frac{\abs v^2}{s^2}+\frac{\alpha^2\abs v^2}{s^3}\right)\dd t.
\]
The cross-variation between the martingale parts of $vv^T$ and $s^{-1}$ contributes
$-\alpha(KB+BK)\dd t$.  Consequently the drift of $B$ is
\begin{equation}\label{eq:sol-dB-before-cancel}
 sK^2-(AB+BA)+rB+\alpha^2\frac r sB-\alpha(KB+BK).
\end{equation}
Now $K=G+\alpha B/s$ and $B^2=rB$.  In
\eqref{eq:sol-dB-before-cancel}, the terms linear in $\alpha$ cancel and the three quadratic
terms have coefficients $1+1-2=0$.  This leaves the drift in \eqref{eq:sol-dB}.

Finally subtract \eqref{eq:sol-dB} from the covariance SDE in
\eqref{eq:sol-loc-martingale}.  Since
\[
 -A^2+AB+BA-rB=-(A-B)^2=-R^2,
\]
we obtain \eqref{eq:sol-dR}.  The calculation is first made before a bounded stopping time;
letting the stopping levels increase proves the local semimartingale identities.
\end{proof}

\begin{theorem}[= Theorem~\ref{thm:scalar-riccati}]
\label{thm:sol-scalar-riccati}
There is a scalar local martingale $M$ such that
\[
 \dd r_t=\dd M_t+(S_t-D_t)\dd t.
\]
Moreover $D_t\ge r_t^2$, so $S_t$ is the only positive drift term.
\end{theorem}

\begin{proof}
Take the trace in \eqref{eq:sol-dB}.  Since
$\Tr(AB)=s\delta^TA\delta$, $\Tr(rB)=r^2$, and
$\Tr(G^2)=\norm G_\HS^2$, its drift is
\[
 s\norm G_\HS^2-2s\delta^TA\delta+r^2=S-D.
\]
For coercivity, $B\preceq A$ implies
\[
 s\abs\delta^4=\delta^TB\delta\le\delta^TA\delta.
\]
Multiplying by $2s$ gives
$2s\delta^TA\delta\ge2s^2\abs\delta^4=2r^2$, hence $D\ge r^2$.
\end{proof}

\begin{corollary}[= Corollary~\ref{cor:per-direction}]
\label{cor:sol-per-direction}
For every $T>0$ and every unit vector $\theta$,
\[
 \E\int_0^T\bigl(s_t\abs{G_t\theta}^2+\abs{R_t\theta}^2\bigr)\dd t
 \le\theta^TR_0\theta\le1.
\]
Consequently
\[
 \E\int_0^\infty s_tG_t^2\dd t\preceq R_0\preceq I_n.
\]
\end{corollary}

\begin{proof}
Apply \eqref{eq:sol-dR} to $\theta^TR_t\theta$.  Up to a scalar local martingale its drift is
\[
 -\theta^TR_t^2\theta-s_t\theta^TG_t^2\theta
 =-\abs{R_t\theta}^2-s_t\abs{G_t\theta}^2.
\]
Stop when the martingale and coefficients are bounded, take expectations, and use
$R_{T\wedge\sigma}\succeq0$.  Fatou's lemma as $\sigma\uparrow\infty$ yields the displayed
finite-$T$ inequality.  Monotone convergence first in $T$ and then testing every deterministic
$\theta$ gives the Loewner inequality on $[0,\infty)$.  Finally
$R_0=A_0-B_0\preceq A_0=I_n$.
\end{proof}

\subsection*{4. Tight-window consumption}

For $0<\eta\le1/6$ define the continuous-exit time
\[
 \tau_\eta=\inf\{t\ge0:\abs{p_t-1/2}>\eta\}.
\]

\begin{corollary}[= Corollary~\ref{cor:tight-window-consumption}]
\label{cor:sol-tight-window-consumption}
Suppose $\abs{p_0-1/2}\le\eta/2$ and, for nonnegative constants $C_0,C_1$, a number
$\alpha<1$, and every $T\le T_0$,
\begin{equation}\label{eq:sol-tight-carleson}
 \E\int_0^{T\wedge\tau_\eta}S_t\dd t
 \le C_0T+C_1\E\int_0^{T\wedge\tau_\eta}r_t\dd t
 +\alpha\E\int_0^{T\wedge\tau_\eta}D_t\dd t.
\end{equation}
Then there is a time $T_*>0$, depending only on
$T_0,C_0,C_1,\alpha,\eta$, for which
\[
 \Prob\{\min(p_{T_*},q_{T_*})\ge1/3\}\ge1/2.
\]
Hence Lemma~\ref{lem:sol-survival-implies-kls} gives a boundary lower bound whose constant has
the same dependence.  In particular, universal input constants give a universal bound.
\end{corollary}

\begin{proof}
Let $u(T)=\E r_{T\wedge\tau_\eta}$.  Stop once more at bounded coefficient/martingale levels,
integrate Theorem~\ref{thm:sol-scalar-riccati}, and then remove that auxiliary stopping.  The
nonnegative terms and the regularization convention justify Fatou where equality is not
needed.  Using \eqref{eq:sol-tight-carleson}, $r_0\le1$, and $1-\alpha>0$ gives
\[
 u(T)\le1+C_0T+C_1\E\int_0^{T\wedge\tau_\eta}r_t\dd t
 \le1+C_0T+C_1\int_0^T u(t)\dd t.
\]
The last inequality uses
$\one_{\{t<\tau_\eta\}}r_t\le r_{t\wedge\tau_\eta}$ pointwise.  Gronwall therefore yields
\begin{equation}\label{eq:sol-u-bound}
 u(T)\le(1+C_0T)e^{C_1T}\le C_*\qquad(0\le T\le T_0),
\end{equation}
where $C_*$ depends only on the displayed input constants.  Before exit,
$s_t\ge1/4-\eta^2\ge2/9$, and so
\begin{equation}\label{eq:sol-centroid-window}
 \E\int_0^{T\wedge\tau_\eta}\abs{\delta_t}^2\dd t
 \le\frac92\int_0^T\E[\one_{\{t<\tau_\eta\}}r_t]\dd t
 \le\frac92C_*T.
\end{equation}

Continuity of $p$ and the nested initial window imply that on $\{\tau_\eta\le T\}$,
$\sup_{t\le T\wedge\tau_\eta}\abs{p_t-p_0}\ge\eta/2$.  Doob's $L^2$ inequality and
\eqref{eq:sol-p-qv} give
\[
 \Prob(\tau_\eta\le T)
 \le\frac4{\eta^2}\E[p]_{T\wedge\tau_\eta}
 =\frac4{\eta^2}\E\int_0^{T\wedge\tau_\eta}s_t^2\abs{\delta_t}^2\dd t
 \le\frac{9C_*T}{8\eta^2},
\]
because $s_t\le1/4$.  Choose $T_*\le T_0$ positive so that the last expression is at most
$1/2$.  On $\{\tau_\eta>T_*\}$,
$\min(p_{T_*},q_{T_*})\ge1/2-\eta\ge1/3$, proving the claim.
\end{proof}

\subsection*{5. Pathwise control away from zero}

\begin{lemma}[= Lemma~\ref{lem:pathwise-BL}]
\label{lem:sol-pathwise-BL}
For every $t>0$,
\[
 s_t\norm{K_t}_\HS^2\le\frac{4\lmax(A_t)}t\le\frac4{t^2}.
\]
\end{lemma}

\begin{proof}
For a symmetric matrix $M$ let
\[
 f_M(x)=(x-a_t)^TM(x-a_t)-\Tr(MA_t),
 \qquad g=\frac{\one_E-p_t}{\sqrt{s_t}}.
\]
The conditional covariance decomposition gives the exact identity
\begin{equation}\label{eq:sol-color-quadratic}
 \E_{\mu_t}[g f_M]=\sqrt{s_t}\,\inner{K_t}{M}.
\end{equation}
Since $\E g=0$ and $\E g^2=1$, Cauchy--Schwarz and the Brascamp--Lieb inequality for the
$t$-uniformly log-concave law $\mu_t$ give
\[
 s_t\inner{K_t}{M}^2
 \le\Var_{\mu_t}(f_M)
 \le\frac1t\E_{\mu_t}\abs{\nabla f_M}^2
 =\frac4t\Tr(MA_tM)
 \le\frac{4\lmax(A_t)}t\norm M_\HS^2.
\]
The posterior's Gaussian factor gives all polynomial moments for $t>0$, so the quadratic
$f_M$ belongs to the Brascamp--Lieb form domain; equivalently one may apply the inequality to
smooth truncations and pass by form closure.
If $K_t\ne0$, take $M=K_t/\norm{K_t}_\HS$; otherwise the assertion is immediate.
The second inequality follows from \eqref{eq:sol-BL-cap}.
\end{proof}

\begin{corollary}[= Corollary~\ref{cor:away-from-zero}]
\label{cor:sol-away-from-zero}
There is a numerical $\eta_0>0$ such that, for
$0<\eta\le\eta_0$ and on $\{t<\tau_\eta\}$,
\[
 S_t\le\frac8{t^2}+\frac12D_t.
\]
Thus, for every interval $I\subset[t_0,T_0]$ with $t_0>0$,
\[
 \E\int_{I\cap[0,\tau_\eta]}S_t\dd t
 \le\frac8{t_0^2}\abs I+\frac12
 \E\int_{I\cap[0,\tau_\eta]}D_t\dd t.
\]
\end{corollary}

\begin{proof}
On the tight window, $\abs{q-p}\le2\eta$.  For $\eta\le1/4$ we also have
$s\ge3/16$.  Since $G=K-(q-p)\delta\delta^T$,
\[
 S=s\norm G_\HS^2
 \le2s\norm K_\HS^2+2s(q-p)^2\abs\delta^4
 =2s\norm K_\HS^2+\frac{2(q-p)^2}{s}r^2.
\]
The first term is at most $8/t^2$ by Lemma~\ref{lem:sol-pathwise-BL}; the second is at most
$(128/3)\eta^2D$ because $r^2\le D$.  Taking, for example,
$\eta_0=\min\{1/4,\sqrt3/16\}$ makes this coefficient at most $1/2$.  Integration on an
interval bounded away from zero gives the final display.
\end{proof}

\paragraph{Endpoint audit.}
Every expectation of a stochastic identity above is taken first at a bounded localizing
stopping time.  Finite-horizon estimates survive its removal by Fatou and the nonnegativity of
$R,S,D$ in the places used.  The only infinite-horizon assertion is obtained from the already
proved finite-horizon occupation inequality by monotone convergence.  The survival step uses
the perimeter supermartingale direction, which is valid without compact support.  Thus no
claim depends on silently declaring a local martingale to be uniformly integrable at time
$\infty$.

\end{document}
